\section{Introduction}
\label{sec:intro}
Speech synthesis and voice conversion have made the assessment of speech authenticity an operational problem for
speech technology. A countermeasure (CM) produces a score, but a screening service or a moderation workflow must turn
that score into a decision. Its threshold is chosen using development data and then applied to speech whose speakers,
recording conditions and synthesis systems may differ from those used to choose it. Assessing this decision requires
more than a measure of how well the scores rank genuine and spoofed speech.

The equal error rate (EER) and minimum tandem detection cost are useful measures of discrimination
\citep{kinnunen18_tdcf,wang20_asvspoof2019,liu23_asvspoof2021}. They use evaluation labels to locate a favourable
operating point. A development threshold has no such access to the target labels. Consequently, a low evaluation EER
does not establish that a threshold selected elsewhere will preserve either the false-rejection budget for genuine
speech or the false-acceptance budget for spoofs. \citet{zhou26_eer_hides} demonstrate this distinction for a speech
deepfake detector transferred to In-the-Wild (ITW).

This distinction has a long history in biometric evaluation. Half total error rate (HTER) at a development threshold
is established practice in authentication and presentation-attack detection
\citep{bengio04_hter,chingovska12_lbp,korshunov16_crossdb}. Application-independent speaker-recognition evaluation
separates discrimination from calibration through actual and minimum detection costs
\citep{brummer06_calibration,vanleeuwen07_intro}. We apply that distinction to speech deepfake countermeasures. The
threshold transfer gap (TTG) is the excess target HTER at the source development threshold over the best HTER achievable
by thresholding the same target scores. It gives the threshold-dependent part of the error a separate name and
reports it beside the error that threshold selection cannot remove.

A threshold-transfer audit must also state whose errors it counts and how its operating point was chosen. Genuine
speech falsely rejected as spoofed may interrupt screening or prompt an incorrect authenticity accusation; a spoof
accepted as genuine may evade the same workflow. Their relative frequencies depend on the threshold policy, the
target distribution and the application costs. Furthermore, a development set with almost separated classes can
admit several nearly equivalent thresholds that behave differently after transfer. These are assessment questions,
even when the underlying detector architecture remains unchanged.

% src: PREREGISTRATION.md; PREREGISTRATION_ROUND2.md sections 1, 3, 4, 5.
We study these questions in two separately registered rounds. Round~1 audits five released checkpoints trained on the
ASVspoof~2019 logical-access data (LA19),
keeping their author-defined score and input conventions. Round~2 asks whether the findings extend to current
detectors, retrained AASIST-L models, the same recipe trained and evaluated on silence-trimmed speech, and additional
channel conditions. Its design was frozen after the round-1 findings were known and before any round-2 development
or target score existed. This chronology distinguishes replication on new scores from a claim that the later
hypotheses were formed without knowledge of the earlier audit.

The contribution is an assessment protocol and an empirical audit rather than a new detection architecture.
Table~\ref{tab:card} sets out the reporting protocol: development-only threshold selection, separate false-rejection
and false-acceptance rates, discrimination and transfer losses, development-threshold ambiguity, speaker-aware
uncertainty, and cost-sensitive interpretation. The audit distinguishes registered tests from exploratory
sensitivity analyses and examines whether its observations concern released checkpoints or a retrained recipe.
It also evaluates a development-designed accept/review/reject policy by the errors it automates and the traffic it
sends to review. The accompanying reproducibility package is described in the data-availability statement.

All detector scores in this article use desktop floating-point inference. The separate deployment study concerns
conversion, quantization and device measurements. The outcomes of the original quantization hypotheses are
disclosed here to preserve the registration record; their analyses belong to that separate study.
